\section{Actual cohomology comparison for affine covers}

\begin{lemma}[Original positive cohomology and the ordered affine complex]
\label{lem:affine-pushforward-ordered-vanishing}
\lean{EnvelopeProperty.AffinePushforwardCohomology.subsingleton_H_iff_ordered_zero}
\leanok
Let $X$ be a separated scheme, $M$ a quasicoherent module, $p:X\to S$
an actual morphism, and $(U_i)$ a finite linearly ordered affine open
cover. For $q\ge0$, the actual group $H^{q+1}(X,M)$ is zero if and
only if degree $q+1$ homology of the original ordered base-linear
\v Cech complex is zero.
\end{lemma}
\begin{proof}
Every nonempty finite intersection in this affine cover is affine by
separatedness. Quasicoherent affine acyclicity therefore applies to
every original intersection. The native acyclic-cover comparison
identifies vanishing of actual positive sheaf cohomology with
exactness of the original \v Cech complex in this degree. Its native
base-linear version preserves this exactness. The actual ordered and
unordered base-linear complexes have the native homology isomorphism,
so exactness is equivalent to zero homology of the ordered complex.
This is the general version of the actual comparison used for closed
projective models; no cohomology-vanishing certificate is assumed.
\end{proof}

\begin{theorem}[Actual affine pushforward preserves positive vanishing]
\label{thm:affine-pushforward-cohomology-vanishing}
\uses{lem:affine-pushforward-ordered-vanishing}
\lean{EnvelopeProperty.AffinePushforwardCohomology.subsingleton_H_iff_pushforward}
\leanok
Let $f:X\to Y$ be affine, $Y$ separated, and $M$ quasicoherent.
For an original finite ordered affine open cover of $Y$ and every
$q\ge0$, the actual $H^{q+1}(X,M)$ vanishes if and only if the actual
$H^{q+1}(Y,f_*M)$ vanishes. There is no Noetherian, finite-type,
field, flatness or finiteness assumption on $f$ beyond affineness.
\end{theorem}
\begin{proof}
An affine morphism is separated; composition with the separated
structure morphism of $Y$ makes $X$ separated. Inverse images of the
original affine opens are affine and cover $X$. The actual $f_*M$
is quasicoherent. Apply the preceding comparison to both original
covers. The native ordered base-linear pushforward isomorphism
identifies the two whole complexes, with their original restrictions
and differentials. Its homology isomorphism identifies zero homology
in either direction, giving exactly the asserted vanishing equivalence.
\end{proof}

\begin{theorem}[Original invertible twists under affine pushforward]
\label{thm:affine-pushforward-twisted-cohomology}
\uses{thm:affine-pushforward-cohomology-vanishing,thm:fc-original-scheme-line-projection-formula}
\lean{EnvelopeProperty.AffinePushforwardCohomology.subsingleton_H_twist_iff_pushforward}
\leanok
In the same geometric setting, let $L$ be an actual invertible module
on $Y$. Then the actual positive cohomology of $M\otimes f^*L$ on $X$
vanishes if and only if that of $f_*M\otimes L$ on $Y$ vanishes.
\end{theorem}
\begin{proof}
The actual pullback of $L$ is invertible, so tensoring the original
quasicoherent $M$ with it remains quasicoherent. Apply the preceding
affine comparison to this tensor product. The already proved genuine
line projection formula identifies its original pushforward with
$f_*M\otimes L$. Applying actual sheaf cohomology to this isomorphism
preserves and reflects vanishing. No projection formula is an extra
hypothesis.
\end{proof}

\begin{lemma}[Actual cohomology descends through a split inclusion]
\label{lem:actual-cohomology-retract-vanishing}
\lean{EnvelopeProperty.AffinePushforwardCohomology.subsingleton_H_of_retract}
\leanok
For actual modules $M,N$ on the same scheme, maps $i:M\to N$,
$r:N\to M$ with $r\circ i=1_M$, and every $q\ge0$, vanishing of
actual $H^q(N)$ implies vanishing of actual $H^q(M)$. No
quasicoherence or geometry hypothesis is needed.
\end{lemma}
\begin{proof}
The composite of the module-to-abelian-sheaf functor and actual
$H^q$ sends the retraction identity to the identity of $H^q(M)$.
Since $H^q(N)$ is zero, the first mapped arrow is zero, so this
identity is zero. An abelian group with zero identity is the zero
group. Degree zero is included. This is a general functorial support
lemma; a cyclic-cover application must first construct the retraction
of its actual original structural unit.
\end{proof}
